\documentclass[./unravbook.tex]{subfiles}

\usepackage{parskip}
\setlength{\parskip}{10pt}
\setstretch{1.15}

\begin{document}

\setcounter{chapter}{0}
\chapter{Reconceptualized} %% {{{
    %% end header ----------------------------

Since October 7, 2023, efforts to diminish Israel's standing in the West seem to have reached a clamorous crescendo.\footcite[In this paper, the \enquote{West} is meant to connote western European and other societies whose influence and power expanded transcontinentally due to the successes of their eighteenth and nineteenth century era of imperialisms. In recent decades, these societies have often been characterized as having imposed their geopolitical will on \enquote{eastern} societies that were inadequately prepared to defend themselves against Western military, industrial and cultural aggressions, which enabled their long successful practice of \enquote{Orientalism.} In Edward Said's words, \enquote{what the early Orientalist achieved, and what the non-Orientalist in the West exploited, was a reduced model of the Orient suitable for the prevailing, dominant culture and its theoretical (and hard after the theoretical, the practical) exigencies.}][153]{Said-1978a} Built on narratives of the country's inherent racism and focusing specifically on its construction and enforcement of a system of apartheid, a resulting challenge that now confronts World Jewry --- and not just Israelis --- is how to assess a concomitant (and arguably corresponding) resurgence in what is often described as the world's \enquote{oldest hatred,} antisemitism. Thus, while Israel can claim recent victories in the military conflicts ignited by its most passionate and lethal enemies, it must be acknowledged that on a global basis, its defeat in the \enquote{war of words} \parencite{Qutami-2020} raises larger questions for a singular people whose presence is nearly worldwide. 

It is worth revisiting the question of Israel, racism and apartheid now from a multidimensional perspective because much has changed since the years of the current century's first decade, when it first became clear that Israel was encountering extreme turbulence on the battlefield of ideas with regard to the subject, and concerned organizations in the Jewish world began consciously preparing for the long-term challenge.\footnote{E.g., Michael Herzog, ``Delegitimization: Attitudes Toward Israel and The Jewish People,'' The Jewish People Policy Institute, \url{https://jppi.org.il/uploads/DELEGITIMIZATION-\%20ATTITUDES\%20TOWARD\%20ISRAEL\%20AND\%20THE\%20JEWISH\%20PEOPLE.pdf}. (``The phenomenon of delegitimization and contending with it is a type of asymmetric warfare, but this time conducted on the battlefield of ideas.'')} That earlier era coincided with the waning of the Oslo Accords process, which had brought a more quiescent period in organized Israel hatred during the 1990s, followed by the violence of the ``Second (or Al-Aqsa) Intifada'' (2000-2005), and followed thereafter by Hamas taking full control of Gaza after an unexpected landslide victory in local elections (2006).\footnote{Anna Mier y Teran, \enquote{War In Palestine and Its Impact on Western Countries,} Universidad de Navarra Global Affairs and Strategic Studies, January 20, 2024. \url{https://www.unav.edu/web/global-affairs/war-in-palestine-and-its-impact-on-western-countries}.} The increasing anathema and change in tone by Israel's enemies was alarming then, but it was not necessarily considered entirely unmanageable; that is, accusations of racism, settler-colonialism, ethnic cleansing, apartheid and much else were hardly new, and arguably could be seen as unlikely to jeopardize the long-term trajectory of Israel's economic, societal and cultural progress \parencite[15--26]{Adelman-2008}.  

The violence that has steadily risen against Jewish communities as a result has understandably caused growing concern that an old era has returned but now in its latest incarnation, where the ``Zionist'' label offers a convenient disguise for acting on one's revulsions toward Jews.\footnote{Johannes Due Enstad (2024). ``Facing antisemitism in Europe: individual and country-level predictors of Jews' victimization and fear across twelve countries.'' \textit{Social Forces}, 103(3), 1186--1209. \url{https://doi.org/10.1093/sf/soae091}.}  The now successfully globalized intifada, or at least one that is in a state of effervescent vitality, has roots that can be easily traced to the mid-2010s,\footnote{A call-to-action was made by the International Union of Muslim Scholars (IUMS), a group led since its 2004 inception by Sheikh Yusuf Al-Qaradawi, ``a leading  Islamic jurist whose pronouncements the Muslim Brotherhood consider authoritative.'' K Shideler and D. Daoud, ``Command and Control: The International Union of Muslim Scholars, The Muslim Brotherhood, and The Call for Global Intifada during Operation Protective Edge,'' Center for Security Policy Occasional Paper Series.  \url{https://www.centerforsecuritypolicy.org/wp-content/uploads/2014/08/Command_and_Control_08-21-14_12.50pm.pdf}. The IUMS made the call to both the entirety of the Muslim \textit{Ummah} and ``free peoples of the world,'' urging them to ``rise up in support of Palestine.''} but the bottom line for everyone against the manning of its barricades is what to do to reverse a now rising tide. This has led to increasing urgency for defensive measures to protect soft targets from rising intimidation and violence, but it should also raise fundamental questions about the bona fides of the racism and apartheid narratives. It no longer seems the right question to ask whether they are ``true'' in any empirical sense, but whether --- accepting that they are planted, poisonous, and poised to grow, and must therefore be expected to continue to threaten anyone open to being receptive toward Israel  --- there are ways to counter the threats they pose from a better informed and more insightful perspective. 

More precisely, UNGA committees, representatives and rapporteurs nowadays have seem to have become the willing midwives of those who seek the birth of the era of Israel's violent destruction. They do so while assuming they are applying the same justified tactics, essentially on the same basis, that were deployed against South Africa during the long decades of its apartheid era. Often led by but by no means limited to Islamist Iran, they have socialized, authenticated and normalized broad anti-Israel public sentiment into the fabric of Western societies through an ever expanding universe of claims and controversies relating to apartheid and racial discrimination. As one of six \enquote{principal organs} of the UN,\footnote{The UNGA is one of ``six principal organs of the United Nations, as established by the UN Charter.'' United Nations (n.d). Main Bodies. \url{https://www.un.org/en/about-us/main-bodies}.} the primary way the UNGA has defined and defended this narrative is by allowing hostile member states and their hirelings to use its forums and impact delivery vehicles to castigate Israel for having ``usurped'' on racist grounds the ``ancestral homeland'' of Palestine.\footnote{Kali Robinson (Last modified July 12, 2023), ``What Is U.S. Policy on the Israeli-Palestinian Conflict?,'' Council on Foreign Relations. \url{https://www.cfr.org/backgrounders/what-us-policy-israeli-palestinian-conflict}.} 

Essentially, such claims are merely elaborative echoes of the powerful slanders of Dr. Fayez Sayegh, an Arab diplomat whose career in the UNGA spanned the critical decades of this long process. His corrupt influence, aided by his immense charm and prolific pen, will be considered in some detail in section IV. But over 60 years ago, and during a brief interlude from his retainer on behalf of various Arab states, whether individually or as a collective, he made the case that still animates how the expanding network of Israel's most diehard enemies --- both within the UNGA and outside it --- still thinks \parencite[214]{Sayegh-1965}:

\begin{adjustwidth}{.65cm}{.5cm}
    \begin{singlespace}
    Racism is not an acquired trait of the Zionist settler-state. Nor is it an accidental, passing feature of the Israeli scene. It is congenial, essential and permanent. For it is inherent in the very ideology of Zionism and in the basic motivation for Zionist colonisation and statehood.
    \end{singlespace}
\end{adjustwidth}

The singular impact of this assertion is clear not only from the UNGA's passage of its infamous \enquote{Zionism as Racism} resolution (3379) ten years later (1975),\footnote{Elimination of all forms of racial discrimination: Zionism as racism --- UNGA resolution 3379, on the report of the Third Committee (A/10320), November 11, 1975, Document A/RES/3379 XXX, \url{https://www.un.org/unispal/document/auto-insert-181963/}.} which basically reiterates it and which he both largely authored and tirelessly evangelized, but also the fact that it derives from a tradition of characterizing Israel --- again, using his words --- as having a ``vital link'' with western imperialism \parencite[21]{Sayegh-1965}. Because racism has come to be widely understood over the past half century as both cause and effect of imperialism's success, the view has become entrenched that Israel's capitulation is only the start of a struggle to end a multi-century era in which the world's ``national capitalist economies'' engineered ``the racialization of the peoples of the entire world'' \parencite[89]{Silva-1999}. The passage of Resolution 3379 gave Israel's adversaries the template for solidifying worldwide opposition in the name of an ``enlightened worldview adopted by good people.''\footnote{Deist, J. (2023). \textit{A Strange Liberty: Politics Drops Its Pretenses}. Ludwig von Mises Institute. \url{https://mises.org/online-book/strange-liberty-politics-drops-its-pretenses/15-getting-liberalism-wrong}. (``In this view, liberalism is a way of thinking, a kind of enlightened worldview adopted by good people. Liberalism becomes almost undefinable by its breadth: an interconnected passel of open-minded and tolerant attitudes across a swath of 'social' issues. Cosmopolitanism, reason, tolerance, feminism, antiracism, and a healthy degree of support for egalitarian democracy are at the fore; economics is not the focus.'')} This, in turn, resulted from a broad shift in academic understandings about the nature of race and racism that occurred during the same era that the UN became operational, with a key ingredient being \enquote{the emergence of a new transnational political culture during the post-war era, that also included other solidarity movements, as well as student's, green, peace and women's movements, often conceptualized as \enquote{the new social movements},} coupled with South African anti-apartheid activism then and Israeli anti-apartheid activism now, creating what has become \enquote{a `movement of movements'.}\footnote{Håkan Thörn, \textit{Anti-Apartheid and the Emergence of a Global Civil Society}, Palgrave: London, 2006, 8.} 

\parindent 0pt
\parskip 6pt


\setenotez{
  list-name = {Notes to Chapter \thechapter}
}
% Place this macro at the exact end of every chapter:
{\footnotesize
\printendnotes[paragraph]}

\thispagestyle{empty} % makes a blank page

\end{document}